\section{Bases chosen for an almost commuting family}

The finite-dimensional argument below constructs a basis from the operators
rather than assuming separated eigenvalue labels. It follows the spectral
refinement in \cite{Hastings2006Solving}, with both compression cross terms
retained as explained in \cite{gap:hastings05_basis_refinement}.
The statements here have no lattice or geometric growth assumption.

\begin{lemma}[Eigenvector entries and off-diagonal parts]
    \label{lem:qic_approx_diagonalization_entries}
    \lean{ApproximateDiagonalization.abs_sub_mul_norm_inner_le_norm_commutator,
    ApproximateDiagonalization.norm_inner_le_norm_commutator_div}
    \leanok

    \lean{ApproximateDiagonalization.offDiagonalPart,
    ApproximateDiagonalization.offDiagonalPart_apply,
    ApproximateDiagonalization.norm_offDiagonalPart_le,
    ApproximateDiagonalization.norm_offDiagonalPart_le_norm_commutator_div,
    ApproximateDiagonalization.sum_norm_offDiagonalPart_le,
    ApproximateDiagonalization.norm_apply_sub_inner_smul_le_add_offDiagonalPart}
    \leanok

    \lean{ApproximateDiagonalization.shiftPart,
    ApproximateDiagonalization.shiftPart_apply,
    ApproximateDiagonalization.norm_shiftPart_le,
    ApproximateDiagonalization.offDiagonalPartFamily,
    ApproximateDiagonalization.norm_offDiagonalPartFamily_le,
    ApproximateDiagonalization.norm_offDiagonalPartFamily_le_norm_commutator_div,
    ApproximateDiagonalization.sum_norm_offDiagonalPartFamily_le,
    ApproximateDiagonalization.offDiagonalPartFamily_apply_self,
    ApproximateDiagonalization.norm_apply_sub_inner_smul_le_add_offDiagonalPartFamily}
    \leanok

    Let $A$ be self-adjoint, with $Av=rv$ and $Aw=sw$. Then
    \[
      |r-s|\,|\langle v,Bw\rangle|
      \le \|[A,B]\|\,\|v\|\,\|w\|.
    \]
    In an orthonormal family $(v_a)_{a<n}$, the operator with the same
    off-diagonal entries as $B$ and zero diagonal has norm at most
    $(n-1)\max_{a\ne c}|\langle v_a,Bv_c\rangle|$.
    A two-vector orthonormal family gives the bound by the larger of its
    two off-diagonal entries. These bounds apply termwise to finite sums.
\end{lemma}
\begin{proof}\leanok
    Take the inner product of the commutator with the two eigenvectors.
    Decompose an off-diagonal matrix into its cyclic shifts. The images of
    the basis vectors under each shift are orthogonal, so its norm is at
    most its largest entry. Sum the $n-1$ shifts. Decomposing an image into
    its projection onto the family and its orthogonal component gives the
    corresponding approximation estimate.
\end{proof}

\begin{lemma}[Restriction and recombination]
    \label{lem:qic_approx_diagonalization_refinement}
    \lean{ApproximateDiagonalization.subspaceCompression,
    ApproximateDiagonalization.subspaceExtension,
    ApproximateDiagonalization.coe_subspaceCompression_apply,
    ApproximateDiagonalization.subspaceExtension_coe_apply,
    ApproximateDiagonalization.inner_subspaceCompression_right,
    ApproximateDiagonalization.inner_subspaceCompression_left,
    ApproximateDiagonalization.norm_subspaceCompression_le,
    ApproximateDiagonalization.norm_subspaceExtension,
    ApproximateDiagonalization.subspaceExtension_compression,
    ApproximateDiagonalization.subspaceExtension_mul,
    ApproximateDiagonalization.subspaceExtension_sub,
    ApproximateDiagonalization.isSymmetric_subspaceCompression,
    ApproximateDiagonalization.isSelfAdjoint_subspaceCompression,
    ApproximateDiagonalization.norm_subspaceCompression_commutator,
    ApproximateDiagonalization.exists_orthonormalBasis_orthogonal_sum,
    ApproximateDiagonalization.exists_orthonormalBasis_orthogonal_sum_fin,
    ApproximateDiagonalization.norm_inner_subspace_cross_le,
    ApproximateDiagonalization.norm_inner_le_of_orthogonal_sum}
    \leanok

    For a subspace $K$ of a finite-dimensional Hilbert space, let $P_K$
    be its orthogonal projection. Compression restricts $P_KAP_K$ to $K$;
    extension by zero preserves the operator norm. Compressions of
    self-adjoint operators are self-adjoint. Orthonormal bases of $K$ and
    $K^\perp$ combine into an orthonormal basis of the original space,
    retaining their vectors exactly. Entries between the two summands
    are bounded by the norm of $P_K A(I-P_K)$ times the vector norms.
\end{lemma}
\begin{proof}\leanok
    Compose the inclusion and orthogonal projection. Their norms are at
    most one, and restriction gives the reverse inequality for extension
    by zero. Transport the product of the two bases along the orthogonal
    direct-sum isomorphism.
\end{proof}

\begin{lemma}[Compression cross terms]
    \label{lem:qic_approx_diagonalization_compression}
    \lean{ApproximateDiagonalization.commutator_compression_eq,
    ApproximateDiagonalization.norm_compression_le,
    ApproximateDiagonalization.norm_crossBlock_le_of_center,
    ApproximateDiagonalization.norm_commutator_compression_le,
    ApproximateDiagonalization.le_sqrt_mul_of_le_spread,
    ApproximateDiagonalization.exp_dist_half_mul_le,
    ApproximateDiagonalization.norm_commutator_compression_le_exp,
    ApproximateDiagonalization.norm_commutator_compression_le_kernel}
    \leanok

    If $P^2=P$ and $Q=I-P$, then
    \[
       [PBP,PCP]=P[B,C]P-PBQCP+PCQBP.
    \]
    Suppose $0\le k\le1$ is symmetric and
    $k(i,j)k(j,l)\le k(i,l)$. If $\|P\|\le1$, the original commutators
    are at most $Ck(i,l)$ and both cross blocks are at most
    $D\sqrt{k(i,j)}$, then the compressed commutators are at most
    $(C+2D^2)\sqrt{k(i,l)}$.
\end{lemma}
\begin{proof}\leanok
    Expand the identity and bound both cross products. Use the
    multiplicative triangle inequality and $k\le\sqrt{k}$.
\end{proof}

\begin{lemma}[An actual spectral split]
    \label{lem:qic_approx_diagonalization_spectral_split}
    \lean{ApproximateDiagonalization.exists_spread_div_le_adjacent_gap,
    ApproximateDiagonalization.exists_adjacent_gap_fin,
    ApproximateDiagonalization.exists_antitone_spectral_split,
    ApproximateDiagonalization.basisSplitProjection,
    ApproximateDiagonalization.basisSplitProjection_apply_basis,
    ApproximateDiagonalization.basisSplitProjection_eq_starProjection,
    ApproximateDiagonalization.norm_basisSplitProjection_le,
    ApproximateDiagonalization.isStarProjection_basisSplitProjection,
    ApproximateDiagonalization.basisSplitProjection_compl,
    ApproximateDiagonalization.basisSplitProjection_mul_mul_compl,
    ApproximateDiagonalization.norm_basisSplitProjection_mul_mul_compl_le,
    ApproximateDiagonalization.norm_basisSplitProjection_cross_le_commutator,
    ApproximateDiagonalization.exists_selfAdjoint_spectral_split,
    ApproximateDiagonalization.norm_le_of_diagonal_basis,
    ApproximateDiagonalization.norm_sub_min_eigenvalue_le_spread,
    ApproximateDiagonalization.exists_family_spectral_split}
    \leanok

    Let $(B_i)$ be a finite self-adjoint family on an $n$-dimensional
    complex Hilbert space, $n>1$, with
    $\|[B_i,B_j]\|\le C f(i,j)$, where $C,f\ge0$ and $f$ is symmetric.
    Either every $B_i$ is scalar, or an original member $B_j$ determines
    a proper nonzero orthogonal spectral projection $P$ such that both
    cross blocks of every $B_i$ have norm at most
    $\sqrt{n^3 C}\sqrt{f(i,j)}$.
\end{lemma}
\begin{proof}\leanok
    Choose a member with maximal spectral spread $\rho$. Zero spread
    gives the scalar alternative. Otherwise one adjacent gap is at least
    $\rho/(n-1)$. Split its eigenbasis there. Centering every member at
    its least eigenvalue bounds each cross block by $\rho$; the
    commutator gives a bound proportional to $Cf/\rho$. Combining the
    two estimates removes the inverse spread.
\end{proof}

\begin{definition}[Successive square-root kernels]
    \label{def:qic_approx_diagonalization_kernel}
    \lean{ApproximateDiagonalization.sqrtKernel,
    ApproximateDiagonalization.sqrtKernel_zero,
    ApproximateDiagonalization.sqrtKernel_succ,
    ApproximateDiagonalization.sqrtKernel_nonneg,
    ApproximateDiagonalization.sqrtKernel_le_one,
    ApproximateDiagonalization.sqrtKernel_symm,
    ApproximateDiagonalization.sqrtKernel_succ_sq,
    ApproximateDiagonalization.sqrtKernel_mono,
    ApproximateDiagonalization.sqrtKernel_mul_le,
    ApproximateDiagonalization.listKernelSum,
    ApproximateDiagonalization.listKernelSum_nil,
    ApproximateDiagonalization.listKernelSum_cons,
    ApproximateDiagonalization.listKernelSum_append,
    ApproximateDiagonalization.listKernelSum_nonneg,
    ApproximateDiagonalization.sqrtKernel_comp}
    \leanok

    Put $k_0=k$ and $k_{r+1}(i,j)=\sqrt{k_r(i,j)}$.
    Positivity, the bound by one, symmetry, and the multiplicative
    triangle inequality persist. The finite list sum is
    $\sum_{j\in\mathcal S} k_r(i,j)$, counting repetitions. Restricting
    a kernel along any map commutes with these operations.
\end{definition}
\begin{proof}\leanok
    Induct on $r$, using monotonicity and multiplicativity of the square
    root on nonnegative real numbers. List sums obey concatenation.
\end{proof}

\begin{theorem}[A constructed basis with weighted kernel bounds]
    \label{thm:qic_approx_diagonalization_uniform_basis}
    \lean{ApproximateDiagonalization.recursiveCommutatorCoeff,
    ApproximateDiagonalization.recursiveBasisCoeff,
    ApproximateDiagonalization.recursiveCommutatorCoeff_nonneg,
    ApproximateDiagonalization.recursiveBasisCoeff_nonneg,
    ApproximateDiagonalization.sqrt_recursiveCommutatorCoeff_le,
    ApproximateDiagonalization.recursiveCommutatorCoeff_step,
    ApproximateDiagonalization.exists_orthonormalBasis_uniform_kernel_aux,
    ApproximateDiagonalization.exists_orthonormalBasis_uniform_kernel}
    \leanok

    \lean{ApproximateDiagonalization.smul_inv_smul_eq_of_norm_le,
    ApproximateDiagonalization.norm_commutator_inv_smul_le,
    ApproximateDiagonalization.exists_orthonormalBasis_uniform_weighted_kernel}
    \leanok

    Let $E$ have complex dimension $n$, and let $(B_i)$ be a finite
    self-adjoint family. Suppose $\|B_i\|\le w_i$, $w_i\ge0$,
    $0\le k\le1$, $k$ is symmetric with the multiplicative triangle
    inequality, and $\|[B_i,B_j]\|\le Cw_iw_jk(i,j)$, $C\ge0$.
    There are an orthonormal basis $(v_a)_{a<n}$ and a list
    $\mathcal S$ of at most $n-1$ original indices such that
    \[
       |\langle v_a,B_i v_c\rangle|
       \le A_n(C)w_i\sum_{j\in\mathcal S}k_n(i,j),\qquad a\ne c,
       \quad A_n(C)=\sqrt{n^3C(1+2n^3)^n}.
    \]
\end{theorem}
\begin{proof}\leanok
    Divide by the norm envelopes; a zero envelope gives a zero operator.
    Induct on the dimension, applying the actual spectral split and
    Lemmas~\ref{lem:qic_approx_diagonalization_refinement}
    and~\ref{lem:qic_approx_diagonalization_compression}.
    Scalar subspaces are terminal irrespective of dimension.
    Both nontrivial children have smaller dimension. The invariant
    $t+\dim E\le n$ controls the depth, while
    $C_t=C(1+2n^3)^t$ controls both retained cross terms.
    Join the two lists with the new selected index and restore the
    envelopes. No inverse-spacing or conditioning bound is assumed.
\end{proof}
